\chapter{Nonparametric Statistics}

\section{Introduction}

Most procedures in Chapters~2--10 assume the data come from a \emph{normal} population. These are \emph{parametric} methods: they rest on one parametric family.

The $t$- and $F$-procedures tolerate slight non-normality. When the population is clearly non-normal, or the data are ranks or categories, use \emph{nonparametric} (\emph{distribution-free}) methods. They:
\begin{itemize}
  \item assume little beyond a continuous distribution;
  \item work with signs and ranks;
  \item keep the stated significance level for many different distributions.
\end{itemize}

\begin{itemize}
  \item \textbf{Advantages:} fewer assumptions; usable for ranked or categorical data; quick to compute; robust to outliers.
  \item \textbf{Disadvantage:} they do not use all the information in the sample. When the population really is normal, they are less efficient than the parametric test and need somewhat larger samples for the same power.
\end{itemize}
Rule of thumb: if the parametric assumptions are justified, use the parametric test. If they are doubtful or cannot be checked, use the nonparametric one.

Many nonparametric methods analyse \emph{ranks}: the positions $1,2,\dots,N$ in the ordered sample.

\section{Sign Test}

\subsection{Description of the Test}

The \emph{sign test} tests the \emph{median} $\tilde\mu$ of a continuous distribution, the value with $\Prob(X\le\tilde\mu)=\Prob(X\ge\tilde\mu)=0.5$.
\begin{itemize}
  \item For a symmetric distribution (such as the normal) the median equals the mean. So the test can replace the $t$-test.
  \item It is valid for \emph{any} continuous distribution.
\end{itemize}

\textbf{Idea.} To test $H_0:\tilde\mu=\tilde\mu_0$, form the differences $X_i-\tilde\mu_0$. Under $H_0$, each difference is equally likely to be positive or negative. So the number of plus signs $R^+$ is binomial with $n$ trials and $p=\tfrac12$.

\begin{definition}[Sign test and its $P$-value]
Let $r^+$ be the number of positive differences and $r^-$ the number of negative ones. Differences equal to zero (\emph{ties}) are discarded, and $n$ is reduced, so $r^++r^-=n$. For $R\sim\mathrm{Bin}(n,\tfrac12)$:
\begin{itemize}
  \item $H_1:\tilde\mu<\tilde\mu_0$: $P=\Prob(R\le r^+)$.
  \item $H_1:\tilde\mu>\tilde\mu_0$: $P=\Prob(R\ge r^+)$.
  \item $H_1:\tilde\mu\ne\tilde\mu_0$: $P=2\Prob(R\ge r^+)$ if $r^+>n/2$, and $P=2\Prob(R\le r^+)$ if $r^+<n/2$.
\end{itemize}
Reject $H_0$ if $P\le\alpha$. Here $\Prob(R\ge r)=\sum_{k=r}^n\binom nk(0.5)^n$.
\end{definition}

\begin{example}\label{ex:sign}
A claim says the median breaking strength of a thread is $50$. Eleven non-tied measurements: $52,48,55,51,47,56,53,49,58,54,57$. (One measurement equal to $50$ was discarded.) Test $H_0:\tilde\mu=50$ against $H_1:\tilde\mu>50$ at $\alpha=0.05$.
\begin{enumerate}[label=\arabic*.]
  \item Parameter: the median breaking strength.
  \item $H_0:\tilde\mu=50$.
  \item $H_1:\tilde\mu>50$.
  \item $\alpha=0.05$.
  \item Statistic: the number of plus differences, $r^+$.
  \item Reject $H_0$ if $P\le0.05$.
  \item The signs are $+,-,+,+,-,+,+,-,+,+,+$, so $r^+=8$ and $n=11$. $P=\Prob(R\ge8)=(165+55+11+1)/2048=0.113$.
  \item Since $0.113>0.05$, do not reject $H_0$. There is no significant evidence that the median exceeds $50$.
\end{enumerate}
\end{example}

\subsubsection*{Critical values}

Tables avoid the binomial sum. Let $r=\min(r^+,r^-)$.
\begin{itemize}
  \item \textbf{Two-sided.} Reject $H_0$ if $r\le r^*_\alpha$. Here $r^*_\alpha$ is the largest value whose tail probability does not exceed $\alpha/2$ in each direction. The actual level is as close to $\alpha$ as the discrete distribution allows, without exceeding it.
  \item \textbf{One-sided.} Put all of $\alpha$ in one tail. For $H_1:\tilde\mu>\tilde\mu_0$, use $r^-$ (reject when the number of minus signs is small). For $H_1:\tilde\mu<\tilde\mu_0$, use $r^+$. Reject if that count is $\le r^*_{2\alpha}$ from the two-sided table. A two-sided level $2\alpha$ is a one-sided level $\alpha$.
\end{itemize}
Do not use the two-sided critical value at the same $\alpha$ for a one-sided test.

\begin{center}
\small
\begin{tabular}{c|cccccccccccccccc}
\toprule
$n$ & 6 & 7 & 8 & 9 & 10 & 11 & 12 & 13 & 14 & 15 & 16 & 17 & 18 & 19 & 20\\
\midrule
$r^*_{0.05}$ & 0 & 0 & 0 & 1 & 1 & 1 & 2 & 2 & 2 & 3 & 3 & 4 & 4 & 4 & 5\\
$r^*_{0.01}$ & -- & -- & 0 & 0 & 0 & 0 & 1 & 1 & 1 & 2 & 2 & 2 & 3 & 3 & 3\\
\bottomrule
\end{tabular}
\end{center}

\subsubsection*{The normal approximation}

For $n$ at least about $10$, $R$ is approximately normal with mean $0.5n$ and variance $0.25n$. So
\[ Z_0=\frac{R^+-0.5n}{0.5\sqrt n}\ \text{ is approximately }\Normal(0,1)\text{ under }H_0. \]
Reject if $\abs{z_0}>z_{\alpha/2}$ (two-sided), $z_0>z_\alpha$ ($H_1:\tilde\mu>\tilde\mu_0$) or $z_0<-z_\alpha$ ($H_1:\tilde\mu<\tilde\mu_0$).

For Example~\ref{ex:sign}: $z_0=(8-5.5)/(0.5\sqrt{11})=1.51$ and $P=0.066$. Same decision.

\begin{pitfall}{ties}
A difference of exactly zero has no sign. The population is assumed continuous, so ties should not occur. If rounding creates them, set them aside, reduce $n$, and apply the test to the rest.
\end{pitfall}

\begin{example}[Propellant shear strength]
Test $H_0:\tilde\mu=2000$ psi against $H_1:\tilde\mu\ne2000$ at $\alpha=0.05$, with $n=20$ motors. The data give $r^+=14$ and $r^-=6$.
\begin{itemize}
  \item \textbf{Exact.} Since $r^+>n/2=10$, $P=2\Prob(R\ge14)=2(0.0577)=0.1153>0.05$. Do not reject.
  \item \textbf{Table.} $r=\min(14,6)=6$. For $n=20$ and $\alpha=0.05$ (two-sided), $r^*=5$. Since $6>5$, do not reject.
  \item \textbf{Normal approximation.} $z_0=(14-10)/(0.5\sqrt{20})=1.789<z_{0.025}=1.96$. Do not reject.
\end{itemize}
All three agree.
\end{example}

\subsection{Sign Test for Paired Samples}

For pairs $(X_{1j},X_{2j})$ from two continuous populations, form $D_j=X_{1j}-X_{2j}$. Apply the sign test to the $n$ differences to test $H_0:\tilde\mu_D=0$ (the two populations have a common median).

\begin{example}
Ten workers are timed (minutes) on a task before and after training. The differences (before $-$ after) have $8$ positive and $2$ negative values.
\begin{itemize}
  \item Two-sided test: $r=\min(8,2)=2$. For $n=10$, $r^*_{0.05}=1$. Since $2>1$, do not reject $H_0$.
  \item Exact $P$-value: $2\Prob(R\le2)=2(1+10+45)/1024=0.109$.
\end{itemize}
\end{example}

\begin{example}[Fuel metering devices]
Twelve cars are each tested with two metering devices. Apply the test to the $12$ paired differences at $\alpha=0.05$. The table gives $r^*=2$. The data give $r=4>2$. So the devices are not shown to differ in median mileage.
\end{example}

\subsection{Type II Error for the Sign Test}

The sign test holds $\alpha$ for any continuous distribution. But $\beta$ depends on the shape of the true distribution. $\beta$ is a binomial probability whose success probability $p$ is the area of the true distribution on the relevant side of $\tilde\mu_0$.

\begin{example}
Take a normal population with $\sigma=1$. Test $H_0:\tilde\mu=0$ (two-sided, $n=12$, $r^*_{0.05}=2$). Suppose the true median is $1.5$.
\begin{itemize}
  \item A negative difference has probability $p=\Prob(X\le0)=\Phi(-1.5)=0.0668$.
  \item The test fails to reject when $R^-\ge3$. So
  \[ \beta=1-\sum_{x=0}^2\binom{12}x(0.0668)^x(0.9332)^{12-x}=1-(0.436+0.375+0.148)=0.042. \]
\end{itemize}
For a non-normal population (say exponential) with the same shift, $p$ and $\beta$ would be very different.
\end{example}

\subsection{Comparison to the $t$-Test}

\begin{itemize}
  \item For a normal population, the $t$-test has the smallest $\beta$ among all tests of the same size. It beats the sign test. The sign test's large-sample efficiency relative to the $t$-test is $2/\pi\approx0.64$.
  \item For symmetric non-normal populations, the $t$-test is still usually better unless the tails are very heavy.
  \item So use the sign test mainly as a quick test about a median, or when only signs are available. The Wilcoxon signed-rank test, which also uses magnitudes, is preferable.
\end{itemize}

\section{Wilcoxon Signed-Rank Test}

The sign test uses only the direction of each difference. The \emph{Wilcoxon signed-rank test} also uses the magnitudes. It applies to \emph{symmetric} continuous distributions, where mean $=$ median. So it tests the mean, $\mu=\mu_0$.

\subsection{Description of the Test}

To test $H_0:\mu=\mu_0$:
\begin{enumerate}
  \item Compute $X_i-\mu_0$.
  \item Rank the absolute differences $\abs{X_i-\mu_0}$ from $1$ (smallest) to $n$.
  \item Attach to each rank the sign of its difference.
\end{enumerate}

\begin{definition}[Signed-rank statistics]
\begin{itemize}
  \item $W^+$: the sum of the ranks with positive signs.
  \item $W^-$: the absolute value of the sum of the ranks with negative signs.
  \item $W=\min(W^+,W^-)$.
\end{itemize}
Since $W^++W^-=n(n+1)/2$, one determines the other.
\end{definition}

Under $H_0$, each rank is equally likely to carry $+$ or $-$. So $W^+$ and $W^-$ should be about equal. The test rejects when they differ a lot.
\begin{itemize}
  \item $H_1:\mu\ne\mu_0$: reject if $w\le w^*$ (table).
  \item $H_1:\mu>\mu_0$: reject if $w^-\le w^*$.
  \item $H_1:\mu<\mu_0$: reject if $w^+\le w^*$. (For one-sided tests use the one-sided significance levels.)
  \item Tied absolute differences get the average of the ranks they would have shared. Zero differences are discarded.
\end{itemize}

\begin{center}
\small
\begin{tabular}{c|cccccccccc}
\toprule
$n$ & 6 & 7 & 8 & 9 & 10 & 11 & 12 & 13 & 14 & 15\\
\midrule
$w^*_{0.05}$ & 0 & 2 & 3 & 5 & 8 & 10 & 13 & 17 & 21 & 25\\
$w^*_{0.10}$ & 2 & 3 & 5 & 8 & 10 & 13 & 17 & 21 & 25 & 30\\
\bottomrule
\end{tabular}
\end{center}
The table gives two-sided levels $0.05$ and $0.10$. One-sided levels are half of these ($0.025$ and $0.05$). So for a one-sided test at $\alpha=0.05$, use the column labelled two-sided $0.10$.

\begin{example}\label{ex:wilcoxon}
Eight pairs give the differences $d=1.2,\ -0.4,\ 2.0,\ 0.8,\ -0.6,\ 1.5,\ 0.9,\ 2.2$. Test $H_0:\mu_D=0$ against $H_1:\mu_D\ne0$ at $\alpha=0.05$.
\begin{itemize}
  \item Absolute values in order: $0.4,0.6,0.8,0.9,1.2,1.5,2.0,2.2$, with ranks $1,\dots,8$.
  \item The negative differences are $-0.4$ (rank $1$) and $-0.6$ (rank $2$). So $w^-=3$, $w^+=36-3=33$, and $w=3$.
  \item For $n=8$, $w^*_{0.05}=3$. Since $w=3\le3$, reject $H_0$. The median difference is not zero at the $5\%$ level.
\end{itemize}
\end{example}

\begin{example}[Textbook examples]
\begin{itemize}
  \item Propellant data ($n=20$, symmetric continuous distribution): $w^+=150$, $w^-=60$, so $w=60$. The two-sided critical value at $\alpha=0.05$ is $w^*=52$. Since $60>52$, do not reject. Same conclusion as the sign test.
  \item Paired fuel-metering data ($n=12$): $w=22.5$ against $w^*=13$. Do not reject.
\end{itemize}
\end{example}

\subsection{Large-Sample Approximation}

If $n\ge20$, $W^+$ (or $W^-$) is approximately normal with
\[ \mu_{W^+}=\frac{n(n+1)}4,\qquad\sigma^2_{W^+}=\frac{n(n+1)(2n+1)}{24}. \]
So $Z_0=\dfrac{W^+-n(n+1)/4}{\sqrt{n(n+1)(2n+1)/24}}$ is approximately standard normal under $H_0$. Use the usual critical regions.

\subsection{Paired Observations}

For pairs from two continuous populations that differ only in their means, rank $\abs{D_j}$ where $D_j=X_{1j}-X_{2j}$, attach signs, and proceed as above. This tests $H_0:\mu_1=\mu_2$ (equivalently $\mu_D=0$).
\begin{itemize}
  \item The \emph{difference} must be symmetric. $X_1$ and $X_2$ themselves need not be.
  \item Ties receive average ranks.
\end{itemize}

\subsection{Comparison to the $t$-Test}

For non-normal data, comparing $\beta$ is hard. So we use the \emph{asymptotic relative efficiency} (ARE): the limiting ratio of the sample sizes needed for the same error probabilities.
\begin{itemize}
  \item Normal populations: ARE of the signed-rank test against the $t$-test is about $0.95$.
  \item Non-normal populations: at least $0.86$, and often above $1$. (Heavy tails hurt the $t$-test because the sample mean is unstable.)
\end{itemize}
So the signed-rank test is never much worse than the $t$-test, and can be much better.

\section{Wilcoxon Rank-Sum Test}

Two independent random samples, of sizes $n_1\le n_2$, come from continuous populations with means $\mu_1,\mu_2$. The distributions are assumed to have the same shape and spread, differing at most in location. Test $H_0:\mu_1=\mu_2$.

\subsection{Description of the Test}

\begin{definition}[Rank-sum statistics]
Arrange all $n_1+n_2$ observations in ascending order and rank them $1,\dots,n_1+n_2$ (ties get average ranks). Let $W_1$ be the sum of the ranks in the \emph{smaller} sample (sample~1) and $W_2$ the sum in the other. Then
\[ W_1+W_2=\frac{(n_1+n_2)(n_1+n_2+1)}2. \]
\end{definition}

If the means are equal, the ranks of the two samples are well mixed, so $W_1$ and $W_2$ are in proportion to the sample sizes.
\begin{itemize}
  \item $H_1:\mu_1\ne\mu_2$: reject if $w_1$ or $w_2$ is at most the tabulated value $w_\alpha$ for $(n_1,n_2)$.
  \item $H_1:\mu_1>\mu_2$: reject if $w_2\le w_\alpha$.
  \item $H_1:\mu_1<\mu_2$: reject if $w_1\le w_\alpha$. (Tabulated one-sided values correspond to $\alpha=0.025$ and $0.005$.)
  \item Examples: $w_{0.05}=78$ for $n_1=n_2=10$. $w_{0.05}=18$ for $n_1=5$, $n_2=6$.
\end{itemize}

\begin{example}\label{ex:ranksum}
Machine~1: $14,18,12,20,17$ ($n_1=5$). Machine~2: $22,25,19,27,24,21$ ($n_2=6$). Test $H_0:\mu_1=\mu_2$ against $H_1:\mu_1\ne\mu_2$ at $\alpha=0.05$.
\begin{itemize}
  \item Ordered, with ranks: $12(1),14(2),17(3),18(4),19(5),20(6),21(7),22(8),24(9),25(10),27(11)$.
  \item Machine~1 has ranks $1,2,4,6,3$, so $w_1=16$ and $w_2=66-16=50$.
  \item The critical value for $(5,6)$ is $18$. Since $16\le18$, reject $H_0$. The machines differ in mean.
\end{itemize}
\end{example}

\begin{example}[Aluminum-lithium alloys]
Ten specimens of a traditional alloy and ten of a new alloy are tested for axial stress. Test $H_0:\mu_1=\mu_2$ at $\alpha=0.05$. For $n_1=n_2=10$, $w_{0.05}=78$ (two-sided). Neither $w_1$ nor $w_2$ is $\le78$, so do not reject. There is no evidence the alloys differ in mean axial stress.
\end{example}

\subsection{Large-Sample Approximation}

When both $n_1,n_2$ exceed $8$, $W_1$ is approximately normal with
\[ \mu_{W_1}=\frac{n_1(n_1+n_2+1)}2,\qquad\sigma^2_{W_1}=\frac{n_1n_2(n_1+n_2+1)}{12}. \]
So $Z_0=(W_1-\mu_{W_1})/\sigma_{W_1}$ is approximately $\Normal(0,1)$ under $H_0$. The critical regions are $\abs{z_0}>z_{\alpha/2}$, $z_0>z_\alpha$ or $z_0<-z_\alpha$, as appropriate.

For Example~\ref{ex:ranksum}: $\mu=30$, $\sigma=\sqrt{30}$, $z_0=(16-30)/5.48=-2.56$, $P=0.011$. (Here $n_1=5$ is below the recommended size, so this is only a rough check.)

\subsection{Comparison to the $t$-Test}

As in the one-sample case:
\begin{itemize}
  \item With normal populations, the rank-sum test is about $95\%$ as efficient as the two-sample $t$-test in large samples.
  \item Whatever the distributions, it is at least $86\%$ as efficient.
  \item It is usually more efficient when the tails are heavier than normal.
\end{itemize}

\section{Nonparametric Methods in the Analysis of Variance}

\subsection{Kruskal--Wallis Test}

The single-factor model of Chapter~7 is $Y_{ij}=\mu+\tau_i+\varepsilon_{ij}$. The \emph{Kruskal--Wallis test} is a nonparametric alternative to its $F$-test. It requires only that the $\varepsilon_{ij}$ have the same continuous distribution for all $a$ treatments.

\textbf{Setup.} Let $N=\sum n_i$. Rank all $N$ observations from $1$ to $N$ (average ranks for ties). Let $R_{i\cdot}$ be the total and $\bar{R}_{i\cdot}$ the average of the ranks in treatment $i$. Under $H_0:\mu_1=\dots=\mu_a$, the expected rank of each observation is $(N+1)/2$.

\begin{theorem}[Kruskal--Wallis statistic]\label{thm:kw}
\[ H=\frac{12}{N(N+1)}\sum_{i=1}^an_i\Big(\bar{R}_{i\cdot}-\frac{N+1}2\Big)^2=\frac{12}{N(N+1)}\sum_{i=1}^a\frac{R_{i\cdot}^2}{n_i}-3(N+1). \]
When $H_0$ is true, and either $a=3$ with all $n_i\ge6$, or $a>3$ with all $n_i\ge5$, $H$ is approximately $\chi^2_{a-1}$. Reject $H_0$ if $h>\chi^2_{\alpha,\,a-1}$.

With many ties, use the corrected statistic
\[ H'=\frac1{S^2}\left[\sum_i\frac{R_{i\cdot}^2}{n_i}-\frac{N(N+1)^2}4\right],\qquad S^2=\frac1{N-1}\left[\sum_{i,j}R_{ij}^2-\frac{N(N+1)^2}4\right]. \]
$S^2$ is the variance of the ranks.
\end{theorem}

\begin{example}\label{ex:kw}
Three treatments, $n_i=5$ each: A: $12,15,11,14,13$. B: $18,20,17,19,16$. C: $14,16,15,13,17$.
\begin{itemize}
  \item Ranking all $15$ values (average ranks for ties) gives rank totals $R_{A\cdot}=19.5$, $R_{B\cdot}=63.0$, $R_{C\cdot}=38.0$. They sum to $120=15\cdot16/2$.
  \item $H=\dfrac{12}{15\cdot16}\cdot\dfrac{19.5^2+63^2+38^2}5-3(16)=0.05(1158.65)-48=9.93$.
  \item With ties: $S^2=(1237.5-960)/14=19.82$ and $H'=(1158.65-960)/19.82=10.02$.
  \item With $a-1=2$ degrees of freedom, $\chi^2_{0.05,2}=5.99$. Reject $H_0$ ($P\approx0.007$). The treatments do not all have the same location.
\end{itemize}
\end{example}

\subsection{Rank Transformation}

A general trick: \emph{replace each observation by its rank} (over all the data, average ranks for ties). Then run the usual normal-theory procedure on the ranks.
\begin{itemize}
  \item For the single-factor design, the $F$-test of Chapter~7 on the ranks is the \emph{rank transformation} test. $H$ is a monotone function of that $F$ statistic, so the two are closely related. It is approximately valid for large samples whatever the distribution.
  \item The same idea gives nonparametric versions of other $F$-tests (blocked designs, regression). A $t$-test on the ranks approximately reproduces the rank-sum test.
\end{itemize}

\textbf{Textbook advice.} If you worry about normality or outliers, run the usual ANOVA on both the original data and the ranks.
\begin{itemize}
  \item If they agree, the ANOVA assumptions are probably fine.
  \item If they disagree, prefer the rank-based result. It is less distorted by non-normality and wild values. Then look into transformations and possible outliers.
\end{itemize}
In the textbook example, five cotton levels are compared at $\alpha=0.01$. There are many ties, so the tie-corrected statistic is used. The critical value is $\chi^2_{0.01,4}=13.28$. Kruskal--Wallis agrees with the ANOVA $F$-test in rejecting $H_0$.

\begin{summarybox}
\begin{itemize}
  \item Nonparametric tests need little beyond continuity and work with signs and ranks. They are slightly less efficient than parametric tests under normality, and often more powerful when tails are heavy.
  \item Sign test: $R^+\sim\mathrm{Bin}(n,\tfrac12)$ under $H_0:\tilde\mu=\tilde\mu_0$. Use exact $P$-values, the table with $r=\min(r^+,r^-)$, or $z_0=(r^+-0.5n)/(0.5\sqrt n)$. Discard ties. For pairs, apply it to the differences. $\beta$ depends on the distribution. The $t$-test is better under normality.
  \item Wilcoxon signed-rank: rank $\abs{X_i-\mu_0}$; $W^+$, $W^-$, $W=\min$; reject if $w\le w^*$. Large $n$: mean $n(n+1)/4$, variance $n(n+1)(2n+1)/24$. ARE $\approx0.95$ against $t$.
  \item Wilcoxon rank-sum: ranks in the combined sample; $W_1$ (smaller sample) and $W_2=(n_1+n_2)(n_1+n_2+1)/2-W_1$. Mean $n_1(n_1+n_2+1)/2$, variance $n_1n_2(n_1+n_2+1)/12$.
  \item Kruskal--Wallis: $H=\frac{12}{N(N+1)}\sum R_{i\cdot}^2/n_i-3(N+1)\approx\chi^2_{a-1}$. Use the tie-corrected $H'$ when there are many ties. The rank transformation turns any ANOVA into a nonparametric test.
\end{itemize}
\end{summarybox}

\section{Exercises}

\begin{exercise}
The following are measurements of the lifetime (hours) of $12$ bulbs: $945,\ 1010,\ 930,\ 1055,\ 990,\ 1100,\ 970,\ 1040,\ 1010,\ 960,\ 1075,\ 985$. Use the sign test to test $H_0:\tilde\mu=1000$ against $H_1:\tilde\mu\ne1000$ at $\alpha=0.05$, find the exact $P$-value and compare with the normal approximation.
\end{exercise}

\begin{exercise}
Eleven pairs of measurements give the paired differences $0.4,\ 0.9,\ -0.2,\ 1.1,\ 0.6,\ 0.3,\ 1.4,\ -0.1,\ 0.8,\ 0.5,\ 1.0$. Test $H_0:\mu_D=0$ with (a) the sign test and (b) the Wilcoxon signed-rank test at $\alpha=0.05$ and compare the conclusions.
\end{exercise}

\begin{exercise}
Compute $\beta$ for the sign test with $n=10$ ($r^*_{0.05}=1$) against a normal alternative whose median is $1.0$ standard deviation above the hypothesised median, assuming the test rejects only for small $r^-$. Compare qualitatively with the $t$-test.
\end{exercise}

\begin{exercise}
Two methods are compared: method~1: $31,\ 28,\ 35,\ 33,\ 29,\ 36,\ 30$ and method~2: $38,\ 41,\ 34,\ 43,\ 40,\ 37,\ 39,\ 42$. Use the Wilcoxon rank-sum test at $\alpha=0.05$ (use the normal approximation) to test for a difference in means, and compare with the two-sample $t$-test.
\end{exercise}

\begin{exercise}
Show that under $H_0$ the Wilcoxon rank-sum statistic $W_1$ has mean $n_1(n_1+n_2+1)/2$. (Hint: every observation has the same expected rank $(n_1+n_2+1)/2$.)
\end{exercise}

\begin{exercise}
Four treatments with five observations each have rank totals $R_{1\cdot}=30,\ R_{2\cdot}=45,\ R_{3\cdot}=75,\ R_{4\cdot}=60$. Compute $H$ and test equality of locations at $\alpha=0.05$.
\end{exercise}

\begin{exercise}
Verify that $W^++W^-=n(n+1)/2$ and that, under $H_0$, $\Exp[W^+]=n(n+1)/4$ (each rank enters $W^+$ with probability $1/2$), and compute the exact distribution of $W^+$ for $n=4$.
\end{exercise}

\begin{exercise}
Explain when you would prefer (a) the sign test, (b) the signed-rank test and (c) the $t$-test for a single sample of $n=15$ measurements, in terms of what is known about the population.
\end{exercise}
